% Source: https://artofproblemsolving.com/community/c3h1765191_infinite_summation_paradox
% Stable identity: aops-topic-1765191-q01
Consider the series $1-\dfrac{1}{2}+\dfrac{1}{3}-\dfrac{1}{4}+\dfrac{1}{5}-\cdots$.

Argument 1: grouping it as $\left(1-\dfrac{1}{2}\right)+\left(\dfrac{1}{3}-\dfrac{1}{4}\right)+\cdots$, every bracket is positive, so the sum is positive.

Argument 2: every positive integer is uniquely $p\cdot2^k$ with $p$ odd and $k\ge0$. Collecting, for each odd $p$, the terms with denominators $p,2p,4p,\dots$ gives $\dfrac{1}{p}\left(1-\dfrac{1}{2}-\dfrac{1}{4}-\cdots\right)=0$, so the sum is $0$.

Which argument is wrong, and why?
